\documentclass[11pt]{article}
\usepackage[margin=1in]{geometry}
\usepackage{amsmath,amssymb}
\usepackage{booktabs}
\usepackage{graphicx}
\usepackage[numbers]{natbib}
\usepackage{hyperref}

\graphicspath{{./}{../outputs/neural_sde_comparison_outputs/}}

\title{Comparison of Classical, Neural ODE, and Neural SDE Fits for Water Kefir Data}
\author{}
\date{}

\begin{document}
\maketitle

\begin{abstract}
    This report compares three dynamical descriptions of water kefir
    fermentation data: a classical logistic ordinary differential equation, a
    deterministic Neural ODE, and a Neural SDE. The data consist of three
    trials measured at 15 time points, giving 45 scalar observations. The
    deterministic Neural ODE is treated as a population model: the three trials
    are interpreted as repeated
    realizations of the same experiment, so one fitted curve is compared with
    all trial observations. The reported residual criteria are in-sample
    diagnostics; out-of-sample MSE should be computed with a time-respecting
    train/test split.
\end{abstract}

\section{Data and preprocessing}

        Let $y_i^{[j]}$ denote the response observed at time $t_i$ in trial$j$,
    with $i=0,\ldots,T-1$ and $j=1,\ldots,J$. In this experiment $T=15$ and
    $J=3$. The time points run from $0$ to $168$ in the original time units.
    The observed index set is
    $$
    \Omega = \{(i,j): y_i^{[j]}\ \hbox{is observed}\},
    \qquad n = |\Omega| = 45.
    $$

        The neural models were trained on normalized time and normalized
    response:
    $$
        \tau_i = \frac{t_i-t_{\min}}{t_{\max}-t_{\min}},
        \qquad
        z_i^{[j]} = \frac{y_i^{[j]}-\bar y}{s_y},
    $$
    where
    $$
        \bar y = \frac{1}{n}\sum_{(i,j)\in\Omega} y_i^{[j]},
        \qquad
        s_y =
            \left[
                \frac{1}{n}\sum_{(i,j)\in\Omega}(y_i^{[j]}-\bar y)^2
            \right]^{1/2}.
    $$
    Thus $\tau_i\in[0,1]$. All model diagnostics were computed on the original
    response scale; for the neural models this required transforming fitted
    normalized values back to the scale of $y_i^{[j]}$.

\section{Models}

    \subsection{Classical logistic ODE}
            The classical baseline is the Verhulst logistic growth model
        \citep{verhulst1838notice}, fitted here as a shared-parameter nonlinear
        regression model. Unlike the neural models, it is fitted directly to the
        original response values $y_i^{[j]}$. The only normalization used by
        this baseline is the normalized time variable $\tau_i\in[0,1]$. Each
        trial has its own observed initial condition, fixed at the first
        measured value, but all trials share the same growth rate and carrying
        capacity:
        $$
            \frac{d y^{[j]}}{d\tau}
            =
            r y^{[j]}\left(1-\frac{y^{[j]}}{K}\right),
            \qquad
            y^{[j]}(0)=y_0^{[j]}=y^{[j]}(t_0).
        $$
        The exact solution used in the fit is
        $$
            y^{[j]}(\tau)
            =
                \dfrac{K}{
                    1 +
                    \left(
                        \dfrac{K - y_0^{[j]}}{ y_0^{[j]} }
                    \right)
                    \exp(-r\tau)
                }.
        $$
        The unknown parameters of the classical fit are only $r$ and $K$; the
        initial values $y_0^{[j]}$ are taken from the data and are not
        optimized. The fit also requires positive observations, because the
        closed-form logistic curve and the ratio $(K-y_0^{[j]})/y_0^{[j]}$ are
        only meaningful for positive initial values.
        Missing observations, if present, are excluded through the observed
        index set $\Omega$.

            The fitted parameters must satisfy $r>0$ and $K>\max_j y_0^{[j]}$.
        Instead of optimizing constrained parameters directly, the
        implementation optimizes two unconstrained real variables
        $a, b \in\mathbb R$ and maps them to valid physical parameters with a
        $\operatorname{softplus}$ transformation:
        $$
            r = \operatorname{softplus}(a)+10^{-6},
                \qquad
                K = K_{\min} +
                \operatorname{softplus}(b)+10^{-6},
        $$
        where $K_{\min}=\max_j y_0^{[j]}+10^{-3}$ and
        $$
            \operatorname{softplus}(x)=\log(1+\exp(x)).
        $$
        The $\operatorname{softplus}$ function is a smooth differentiable
        approximation to the positive part function. It maps every real number
        to a positive value, behaves approximately like $\exp(x)$ for large
        negative $x$, and behaves approximately like $x$ for large positive
        $x$. This is useful for gradient-based fitting:
        Adam can update the unconstrained variables $a$ and $b$ freely, while
        the reported parameters $r$ and $K$ remain positive and the carrying
        capacity stays above every initial condition. The small constants
        $10^{-6}$ and $10^{-3}$ are numerical safeguards that avoid zero growth
        and avoid a carrying capacity equal to an initial value.

The raw parameters are initialized by applying the inverse softplus map. In the
implementation, this gives an initial growth rate near $3.0$ on the normalized
time scale, and an initial carrying capacity slightly above the largest
observed response. At each training epoch the code applies the softplus
transformation and evaluates the exact logistic solution for all trials and all
observation times.

The fitting criterion is the masked original-scale mean squared error
$$
  \mathcal L_{\mathrm{logistic}}(r,K)
  =
  \frac{1}{n}
  \sum_{(i,j)\in\Omega}
  \left(y_i^{[j]}-\widehat y^{[j]}(\tau_i)\right)^2.
$$
This can also be interpreted as a maximum-likelihood fit after adding a
Gaussian observation-error model on the original response scale:
$$
  y_i^{[j]}
  =
  \widehat y^{[j]}(\tau_i;r,K)+\varepsilon_i^{[j]},
  \qquad
  \varepsilon_i^{[j]}\stackrel{\mathrm{ind}}{\sim}
  \mathcal N(0,\sigma^2).
$$
For fixed $\sigma^2$, maximizing this Gaussian likelihood over $r$ and $K$ is
equivalent to minimizing the sum of squared residuals. If $\sigma^2$ is also
unknown, profiling it out gives
$$
  \widehat\sigma^2(r,K)
  =
  \frac{1}{n}
  \sum_{(i,j)\in\Omega}
  \left(y_i^{[j]}-\widehat y^{[j]}(\tau_i;r,K)\right)^2,
$$
and the maximizing values of $r$ and $K$ are again the nonlinear least-squares
solution. Therefore the implemented classical fit is a profiled Gaussian
maximum-likelihood estimate under an independent additive measurement-error
assumption. It is not a transition likelihood for a stochastic logistic process.

The code backpropagates through the closed-form expression above and applies an
Adam update \citep{kingma2015adam} to $a$ and $b$. No numerical ODE solver is
used for this classical baseline. This model was trained for 3000 Adam epochs
with learning rate $5\times10^{-2}$.
The fitted values were
$$
  \widehat r = 7.705769,
  \qquad
  \widehat K = 47.860203.
$$
Since $\tau=t/168$ for these data, the growth rate in the original time scale
is approximately $\widehat r/168=0.045868$ per original time unit.

\subsection{Neural ODE}

The deterministic neural model is the non-autonomous Neural ODE
\citep{chen2018neuralode}
$$
  \frac{d z}{d\tau} = f_\theta(\tau,z),
  \qquad
  z(0)=z_0.
$$
The fitted solution is a single population curve, not one curve per trial. The
three trials are treated as replicate observations of the same underlying
trajectory. On the original response scale the initial condition is
$$
  y_0 = y(t_0)
  =
  \frac{1}{J}\sum_{j=1}^J y_0^{[j]},
$$
the mean of the first observed values across trials. The corresponding
normalized initial condition is
$$
  z_0 = \frac{y_0-\bar y}{s_y}.
$$
After solving the normalized ODE, the fitted value on the original scale is
$$
  \widehat y_i
  =
  \widehat y(t_i)
  =
  \bar y+s_y\widehat z(\tau_i;\theta).
$$
The same value $\widehat y_i$ is compared with every trial observation
$y_i^{[j]}$ at time $t_i$.

The right-hand side $f_\theta$ is a fully connected neural network with input
$(\tau,z)\in\mathbb R^2$, two hidden layers, 32 hidden units per layer, and
hyperbolic tangent activation:
$$
  h_1 = \tanh(W_1 x+b_1),\qquad
  h_{\ell}=\tanh(W_{\ell}h_{\ell-1}+b_{\ell}),\qquad
  f_\theta(x)=W_o h_2+b_o,
$$
where $x=(\tau,z)^T$. The network has 1185 trainable weights and biases.
This count comes from the three affine layers:
$$
\begin{aligned}
  \hbox{input to hidden 1:}&\quad 2\cdot 32+32=96,\\
  \hbox{hidden 1 to hidden 2:}&\quad 32\cdot 32+32=1056,\\
  \hbox{hidden 2 to output:}&\quad 32\cdot 1+1=33.
\end{aligned}
$$
Thus $96+1056+33=1185$. The architecture was chosen as a compact smooth
function approximator for a scalar time-dependent vector field. The input has
only two coordinates, time and state, and the data set has only 45 scalar
observations, so a very large network would make the comparison dominated by
overparameterization. Two hidden layers with 32 units give substantially more
flexibility than the two-parameter logistic curve while remaining small enough
to train stably with weight decay. The $\tanh$ nonlinearity gives a smooth
right-hand side, which is useful for ODE/SDE integration, and sigmoidal neural
networks are classical universal approximators \citep{cybenko1989approximation}.
The same architecture is used for the Neural ODE and for each Neural SDE
coefficient so that differences between the deterministic and stochastic neural
models are not caused by different network sizes.

The ODE was solved at the observed normalized times using RK4 with step size
$0.01$. The empirical loss was
$$
  \mathcal L_{\mathrm{ODE}}(\theta)
  =
  \frac{1}{n}
  \sum_{(i,j)\in\Omega}
  \left(z_i^{[j]}-\widehat z_i(\theta)\right)^2.
$$
Equivalently, in original response units this loss fits one curve
$\widehat y_i$ to all three replicate measurements at the same time. The
minimizer is therefore encouraged to pass through the center of the replicate
cloud, while the ODE structure enforces a smooth dynamical trajectory.
Training used Adam for 1000 epochs with learning rate $10^{-2}$ and weight
decay $10^{-4}$ \citep{kingma2015adam}.

\subsection{Neural SDE}

The stochastic neural model is an Ito diffusion with neural drift and diffusion
coefficients, following the Neural SDE modeling idea
\citep{tzen2019neuralsde,kidger2021neuralsdegan}:
$$
  dZ_\tau
  =
  \mu_\theta(\tau,Z_\tau)\,d\tau
  +
  \sigma_\phi(\tau,Z_\tau)\,dW_\tau.
$$
Both the drift $\mu_\theta$ and the raw diffusion network use the same
architecture as the Neural ODE: two hidden layers, 32 units per hidden layer,
and $\tanh$ activation. The diffusion is constrained to be positive by
$$
  \sigma_\phi(\tau,z)
  =
  \operatorname{softplus}(g_\phi(\tau,z))+\sigma_{\min},
  \qquad
  \sigma_{\min}=10^{-3}.
$$
The total trainable parameter count is 2370, one copy of the 1185-parameter
network for the drift and one for the diffusion.

The SDE was trained with an Euler-Maruyama Gaussian transition likelihood
\citep{kloeden1992numerical}. For successive observations separated by
$\Delta\tau_i=\tau_{i+1}-\tau_i$, the approximation is
$$
  Z_{i+1}^{[j]}\mid Z_i^{[j]}
  \sim
  \mathcal N\left(
    Z_i^{[j]}+\mu_\theta(\tau_i,Z_i^{[j]})\Delta\tau_i,\,
    \sigma_\phi(\tau_i,Z_i^{[j]})^2\Delta\tau_i
  \right).
$$
With the variance floor used in the implementation, define
$$
  m_i^{[j]}=z_i^{[j]}+\mu_\theta(\tau_i,z_i^{[j]})\Delta\tau_i,
  \qquad
  v_i^{[j]}=\sigma_\phi(\tau_i,z_i^{[j]})^2\Delta\tau_i+10^{-6}.
$$
The minimized negative log-likelihood is
$$
  \mathcal L_{\mathrm{SDE}}(\theta,\phi)
  =
  \frac{1}{|\Omega_{\mathrm{tr}}|}
  \sum_{(i,j)\in\Omega_{\mathrm{tr}}}
  \frac{1}{2}
  \left[
    \log(2\pi v_i^{[j]})
    +
    \frac{(z_{i+1}^{[j]}-m_i^{[j]})^2}{v_i^{[j]}}
  \right],
$$
where $\Omega_{\mathrm{tr}}$ is the set of adjacent observed transitions.
Training used Adam for 1000 epochs with learning rate $5\times10^{-3}$ and
weight decay $10^{-4}$ \citep{kingma2015adam}. After training, 1000 Monte Carlo
paths were simulated to estimate predictive means and central 90 percent
intervals.

\section{Comparison criteria}

All model comparison metrics were computed on the original response scale using
the fitted mean trajectory for each model. For models that produce one curve per
trial, write this fitted value as $\widehat y_i^{[j]}$. For the single-curve
Neural ODE, $\widehat y_i^{[j]}=\widehat y_i$ for every trial $j$. Let
$$
  \mathrm{SSE}
  =
  \sum_{(i,j)\in\Omega}(y_i^{[j]}-\widehat y_i^{[j]})^2,
  \qquad
  \mathrm{SST}
  =
  \sum_{(i,j)\in\Omega}(y_i^{[j]}-\bar y)^2.
$$
The root mean squared error is
$$
  \mathrm{RMSE}
  =
  \sqrt{\frac{\mathrm{SSE}}{n}},
$$
and the coefficient of determination is
$$
  R^2
  =
  1-\frac{\mathrm{SSE}}{\mathrm{SST}}.
$$

The numerical table below reports in-sample residual diagnostics: the models
are trained on all observed entries in $\Omega$, and the same entries are used
to compute RMSE and $R^2$. This is appropriate for checking interpolation
quality, but it is not an out-of-sample generalization error. To estimate a
test MSE for this time-series problem, the split should be made by time rather
than by randomly sampling individual observations. Let
$$
  I_{\mathrm{train}}\subset \{0,\ldots,T-1\},
  \qquad
  I_{\mathrm{test}}\subset \{0,\ldots,T-1\},
  \qquad
  I_{\mathrm{train}}\cap I_{\mathrm{test}}=\varnothing,
$$
and define
$$
  \Omega_{\mathrm{train}}
  =
  \{(i,j)\in\Omega:i\in I_{\mathrm{train}}\},
  \qquad
  \Omega_{\mathrm{test}}
  =
  \{(i,j)\in\Omega:i\in I_{\mathrm{test}}\}.
$$
The Neural ODE is then trained by replacing $\Omega$ with
$\Omega_{\mathrm{train}}$ in $\mathcal L_{\mathrm{ODE}}(\theta)$, while the test
MSE is computed as
$$
  \mathrm{MSE}_{\mathrm{test}}
  =
  \frac{1}{|\Omega_{\mathrm{test}}|}
  \sum_{(i,j)\in\Omega_{\mathrm{test}}}
  \left(y_i^{[j]}-\widehat y_i\right)^2.
$$
For these data, a natural forecasting split is to train on early times and test
on the last few time points, for example training through $t=132$ and testing
at $t=144,156,168$. With only 15 time points, an alternative is a rolling
leave-last-time-out protocol, such as training through $t=120$ and testing at
$t=132$, then training through $t=132$ and testing at $t=144$, and so on.

For AIC and BIC \citep{akaike1974newlook,schwarz1978estimating}, residuals were
modeled as independent Gaussian errors with
common variance $\sigma^2$. The maximum likelihood variance estimate is
$$
  \widehat\sigma^2 = \frac{\mathrm{SSE}}{n},
$$
which gives the maximized log-likelihood
$$
  \log \widehat L
  =
  -\frac{n}{2}
  \left[
    \log(2\pi)+1+\log(\widehat\sigma^2)
  \right].
$$
The criteria are
$$
  \mathrm{AIC}=2k-2\log\widehat L,
  \qquad
  \mathrm{BIC}=k\log(n)-2\log\widehat L.
$$
Here $k=p+1$, where $p$ is the number of fitted dynamical parameters and the
additional parameter is the residual variance. Thus $k=3$ for the classical
logistic model, $k=1186$ for the Neural ODE, and $k=2371$ for the Neural SDE.
Initial conditions are not counted as fitted parameters because they are taken
directly from the first row of the data; for the single-curve Neural ODE this
means using $y_0=J^{-1}\sum_{j=1}^J y_0^{[j]}$.

The 90 percent interval coverage is computed only for the Neural SDE, because
only this model produces a predictive distribution rather than a single fitted
mean curve. Let $B=1000$ be the number of simulated SDE paths and let
$Y_i^{(b)[j]}$ denote the simulated original-scale value for path $b$, time
$i$, and trial $j$. For interval level $\gamma=0.90$, set
$\alpha=(1-\gamma)/2=0.05$ and define the pointwise empirical interval
$$
  L_i^{[j]}=Q_{\alpha}\left(\{Y_i^{(b)[j]}\}_{b=1}^B\right),
  \qquad
  U_i^{[j]}=Q_{1-\alpha}\left(\{Y_i^{(b)[j]}\}_{b=1}^B\right),
$$
where $Q_q$ is the empirical $q$-quantile. The reported empirical coverage is
$$
  \widehat C_{0.90}
  =
  \frac{1}{n}
  \sum_{(i,j)\in\Omega}
  \mathbf 1\{L_i^{[j]}\le y_i^{[j]}\le U_i^{[j]}\}.
$$
Thus an observed data point is counted as covered if it lies between the 5
percent and 95 percent Monte Carlo quantiles for the same time and trial. The
reported value $0.9333$ means that 42 of the 45 observed responses fall inside
their corresponding central 90 percent intervals. The mean interval width is
the average of $U_i^{[j]}-L_i^{[j]}$ over the same observed entries. This is an
in-sample empirical calibration diagnostic for the fitted Neural SDE, not an
out-of-sample guarantee.

\section{Numerical results}

\begin{table}[htbp]
\centering
\caption{Model comparison on the original response scale. Lower RMSE, AIC, and
BIC are better; larger $R^2$ is better.}
\begin{tabular}{lrrrrrr}
\toprule
Model & RMSE & $R^2$ & AIC & BIC & $p$ & $k$ \\
\midrule
Classical logistic ODE & 1.4291 & 0.9826 & 165.84 & 171.26 & 2 & 3 \\
Neural ODE             & 1.0281 & 0.9910 & 2502.20 & 4644.90 & 1185 & 1186 \\
Neural SDE             & 2.0312 & 0.9649 & 4933.48 & 9217.08 & 2370 & 2371 \\
\bottomrule
\end{tabular}
\label{tab:metrics}
\end{table}

\begin{table}[htbp]
\centering
\caption{Additional stochastic-model diagnostics.}
\begin{tabular}{lrr}
\toprule
Model & 90 percent interval coverage & Mean interval width \\
\midrule
Neural SDE & 0.9333 & 7.3887 \\
\bottomrule
\end{tabular}
\label{tab:sde-intervals}
\end{table}

The Neural ODE has the best in-sample interpolation accuracy, with the smallest
RMSE and largest $R^2$. The classical logistic ODE is less flexible, but it
still explains approximately 98.3 percent of the observed variance. Because
AIC and BIC penalize the number of fitted parameters, the two-parameter
logistic model is strongly preferred by these information criteria. The Neural
SDE has the largest RMSE among the three fitted mean curves, but its 90 percent
predictive interval coverage is close to the nominal level.

The AIC and BIC values for the neural models should be interpreted as
parameter-penalized residual criteria for the fitted mean curves. They are not
the same as comparing the SDE transition likelihood directly to an exact ODE
likelihood. Their practical role here is to quantify the tradeoff between
in-sample residual accuracy and model complexity.

\begin{figure}[htbp]
\centering
\includegraphics[width=0.96\textwidth]{water_kefir_neural_dynamics_comparison.png}
\caption{Observed trial data, the single Neural ODE population curve, fitted
mean trajectories from the comparison models, and the Neural SDE central
90 percent predictive interval. The Neural ODE is fitted as one curve against
all three trial realizations, while the logistic ODE captures the saturating
growth pattern with a much smaller parameter count.}
\label{fig:fit-comparison}
\end{figure}

\begin{figure}[htbp]
\centering
\includegraphics[width=0.96\textwidth]{training_objectives.png}
\caption{Training objective histories from the saved experiment outputs. The
three objectives are not on the same statistical scale: the classical model
uses original-scale MSE, the Neural ODE uses normalized MSE, and the Neural SDE
uses normalized transition negative log-likelihood. The plots are therefore
optimization diagnostics, not direct model-comparison metrics.}
\label{fig:training-objectives}
\end{figure}

\section{Reproducibility}

The comparison was generated by the \texttt{kefir-sde-compare} command-line
tool, implemented in \path{src/kefir_models/sde_compare.py}. The data file
used in this directory is named \path{data/raw/waterKefirTrialsReferece.csv}.
The following command
reproduces the reported outputs from the project root:
\begin{verbatim}
kefir-sde-compare \
  --data-path data/raw/waterKefirTrialsReferece.csv \
  --output-dir outputs/neural_sde_comparison_outputs \
  --classical-epochs 3000 \
  --classical-learning-rate 5e-2 \
  --ode-epochs 1000 \
  --ode-learning-rate 1e-2 \
  --sde-epochs 1000 \
  --sde-learning-rate 5e-3 \
  --hidden-units 32 \
  --hidden-layers 2 \
  --weight-decay 1e-4 \
  --min-diffusion 1e-3 \
  --ode-method rk4 \
  --sde-paths 1000 \
  --interval-level 0.90 \
  --seed 7 \
  --device auto \
  --plot
\end{verbatim}

The main output files used in this report are:
\begin{itemize}
  \item \path{outputs/neural_sde_comparison_outputs/model_comparison_metrics.csv},
  \item \path{outputs/neural_sde_comparison_outputs/classical_logistic_parameters.csv},
  \item \path{outputs/neural_sde_comparison_outputs/water_kefir_neural_dynamics_comparison.csv},
  \item \path{outputs/neural_sde_comparison_outputs/water_kefir_neural_dynamics_comparison.png}.
\end{itemize}

\bibliographystyle{plainnat}
\bibliography{kefir_model_comparison_references}

\end{document}
